\documentclass[../../main.tex]{subfiles}
\begin{document}
\Needspace{12\baselineskip}
% SOURCE p=185; printed=170; object=ch05-tvs-exercises-title; region=middle
\subsection*{习题}
% SOURCE p=185; printed=170; object=ch05-ex43-seminorm-metrization; number=43; region=middle
\begin{exercise}\label{exercise:ch05-prove-seminorm-metrization}
证明\cref{proposition:ch05-seminorm-metrization}。（第二问仿照\cref{section:ch04-lch-spaces} 中\subcref{exercise:ch04-bounded-function-metrics}{item:ch04-compact-uniform-function-metric} 的方法进行。）
\end{exercise}
% SUPPLEMENT: id=v14-ch05-ex43-solution; parent=ch05-ex43-seminorm-metrization; kind=solution+motivation+hypotheses; origin=independent-derivation
\begin{solution}
\textbf{两问完整补写（重构新增）。}
\begin{enumerate}
\item 先设半范数族分离非零向量。对 $x\ne y$ 选 $\alpha$ 使 $a=p_\alpha(x-y)>0$。若 $z$ 同时落在 $U_{x\alpha,a/3}$ 与 $U_{y\alpha,a/3}$ 中，则 $a\le p_\alpha(x-z)+p_\alpha(z-y)<2a/3$，矛盾。因此这两个开邻域不交，空间 Hausdorff。

反之，若 $v\ne0$ 而所有 $p_\alpha(v)=0$，则任何零点基本邻域都含 $v$，从而任何零点邻域都含 $v$，不能与 $v$ 的邻域分离。因此空间不 Hausdorff，证明第一问的反方向。
\item 对可数非空族编号为 $p_1,p_2,\ldots$；若族有限，末尾补零半范数。定义
\begin{align*}
\rho(x,y)=\sum_{j=1}^{\infty}2^{-j}\min\{1,p_j(x-y)\}.
\end{align*}
级数由 $\sum_j2^{-j}=1$ 控制，故有限非负。对称性来自 $p_j(-u)=p_j(u)$；当 $\rho(x,y)=0$ 时，每项均为零，半范数分离性给 $x=y$。对任意 $a,b\ge0$，$\min(1,a+b)\le\min(1,a)+\min(1,b)$，结合半范数三角不等式逐项相加，得到 $\rho$ 的三角不等式。差在共同平移下不变，所以 $\rho$ 平移不变。

证明它生成原拓扑。给定 $x$ 的基本邻域 $V=\bigcap_{l=1}^m\{y:p_{j_l}(y-x)<\epsilon_l\}$，若 $m>0$，取
\begin{align*}
0<\delta<\min_{1\le l\le m}2^{-j_l}\min(1,\epsilon_l).
\end{align*}
当 $\rho(x,y)<\delta$ 时，每个 $\min(1,p_{j_l}(y-x))<\min(1,\epsilon_l)$，从而 $p_{j_l}(y-x)<\epsilon_l$，所以 $B_\rho(x,\delta)\subseteq V$。空交情形无需限制。

反过来，给 $r>0$，选 $N$ 使 $\sum_{j>N}2^{-j}<r/2$。若 $p_j(y-x)<r/2$ 对 $1\le j\le N$ 成立，则前 $N$ 项之和小于 $r/2$，故 $\rho(x,y)<r$。这样原拓扑的一个基本邻域包含于每个度量球。两种拓扑相同。

若原半范数族为空而空间 Hausdorff，第一问说明空间只能是 $\{0\}$；其唯一的度量与原拓扑相容。因而所有可数情形均已覆盖。
\end{enumerate}
\end{solution}
\begin{remark}
\textbf{动机与逐条件说明（重构新增）。}截断把每种误差压在1以内，可和权重使总距离有限；任意有限组误差仍可从总距离中逐项恢复。共同核为零是正定性的必要条件，可数性使该级数构造可行。平移不变性由差的形式自动保证，并不要求度量对数乘齐次，故不能把这个度量直接当成范数。
\end{remark}
% END SUPPLEMENT: id=v14-ch05-ex43-solution
% SOURCE p=185; printed=170; object=ch05-ex44-sequential-completeness; number=44; region=middle
\begin{exercise}\label{exercise:ch05-sequential-completeness-nets}
若 $\mathcal X$ 为第一可数拓扑向量空间，且其中每个 Cauchy 列都收敛，则其中每个 Cauchy 网都收敛。
\end{exercise}
% SUPPLEMENT: id=v14-ch05-ex44-solution; parent=ch05-ex44-sequential-completeness; kind=solution+motivation+hypotheses; origin=independent-derivation
\begin{solution}
\textbf{不使用共尾子网的完整证明（重构新增）。}取零点的可数开邻域基 $U_1,U_2,\ldots$，以有限交替代可令它递减。加法在 $(0,0)$ 连续，且取负为同胚，所以对每个 $n$ 可取对称开邻域 $V_n$，使 $V_n+V_n\subseteq U_n$；因为 $0\in V_n$，也有 $V_n\subseteq U_n$。

设 $\{x_i\}_{i\in I}$ 是 Cauchy 网。对每个 $n$ 取 $a_n$，使 $i,j\succeq a_n$ 时 $x_i-x_j\in V_n$。递归选 $i_n\succeq a_n,i_{n-1}$（第一步只要求 $i_1\succeq a_1$）。若 $k,l\ge n$，则 $i_k,i_l\succeq a_n$，所以 $x_{i_k}-x_{i_l}\in V_n\subseteq U_n$。由于 $\{U_n\}$ 为邻域基，序列 $x_{i_n}$ 是 Cauchy 列，依假设收敛于某个 $x$。

固定 $n$，选 $k\ge n$ 足够大，使 $x_{i_k}-x\in V_n$。对任意 $i\succeq i_n$，两指标 $i,i_k$ 都在 $a_n$ 之后，故
\begin{align*}
x_i-x=(x_i-x_{i_k})+(x_{i_k}-x)\in V_n+V_n\subseteq U_n.
\end{align*}
这说明原网最终落入每个 $x+U_n$，所以原网收敛于 $x$。
\end{solution}
\begin{remark}
\textbf{动机与条件说明（重构新增）。}第一可数性使需要控制的邻域可以依次列出；网的 Cauchy 性负责把所选序列的极限传回整个网。所选指标不必在 $I$ 中共尾，因此不能声称该序列是原网的子网。Hausdorff 性并未使用：结论只要求存在极限，不要求唯一。逆方向无需第一可数性，因为序列本来就是网。
\end{remark}
% END SUPPLEMENT: id=v14-ch05-ex44-solution
% SOURCE p=185; printed=170; object=ch05-ex45-smooth-real-line; number=45; region=middle
\begin{exercise}\label{exercise:ch05-smooth-real-line-frechet}
$\mathbb R$ 上所有无限次可微函数组成的空间 $C^\infty(\mathbb R)$ 有一个 Fr\'echet 空间拓扑，使 $f_n\to f$ 当且仅当对所有 $k\ge0$，$f_n^{(k)}\to f^{(k)}$ 在紧集上一致收敛。
\end{exercise}
% SUPPLEMENT: id=v14-ch05-ex45-solution; parent=ch05-ex45-smooth-real-line; kind=solution+motivation+hypotheses; origin=independent-derivation
\begin{solution}
\textbf{补写解答（重构新增）。}对 $m\ge1,k\ge0$ 定义 $p_{m,k}(f)=\sup_{|x|\le m}|f^{(k)}(x)|$。它们为有限半范数，族可数，且 $p_{m,0}$ 合起来分离函数，故生成 Hausdorff 局部凸可度量拓扑。每个紧集都包含于某个 $[-m,m]$，所以所有 $p_{m,k}(f_n-f)\to0$ 与题中的逐阶紧一致收敛等价。

设 $f_j$ 关于这些半范数为 Cauchy 列。对每对 $(m,k)$，$f_j^{(k)}$ 在 $[-m,m]$ 上一致趋于连续函数 $g_{m,k}$。在重叠区间上极限相同，故对每个 $k$ 可拼成 $\mathbb R$ 上连续函数 $g_k$。对任意 $s,t\in\mathbb R$，选 $m$ 覆盖二者及其间区间，在
\begin{align*}
f_j^{(k)}(t)-f_j^{(k)}(s)=\int_s^t f_j^{(k+1)}(u)\,\mathrm du
\end{align*}
中取一致极限，得到 $g_k(t)-g_k(s)=\int_s^t g_{k+1}(u)\,\mathrm du$。微积分基本定理给 $g_k'=g_{k+1}$，所以 $g_0\in C^\infty(\mathbb R)$，且每个 $p_{m,k}(f_j-g_0)\to0$。由前题，序列完备即完备，因此这是所求 Fr\'echet 拓扑。
\end{solution}
\begin{remark}
\textbf{动机与条件说明（重构新增）。}两个可数指标分别控制空间范围和导数阶数。紧区间使上确界有限，包含零阶控制保证极限函数可辨别；所有阶的积分相容性保证光滑性保留。此处不要求各阶导数在整个实线上有界，也不要求不同阶共享一个统一界。
\end{remark}
% END SUPPLEMENT: id=v14-ch05-ex45-solution
% SOURCE p=185; printed=170; object=ch05-ex46-initial-seminorm; number=46; region=middle
\begin{exercise}\label{exercise:ch05-initial-seminorm-topology}
若 $\mathcal X$ 为向量空间，$\mathcal Y$ 为赋范线性空间，$\mathcal T$ 为线性映射族 $\{T_\alpha:\mathcal X\to\mathcal Y\}$ 在 $\mathcal X$ 上生成的弱拓扑，$\mathcal T'$ 为半范数族 $\{x\mapsto\|T_\alpha x\|\}$ 定义的拓扑，则 $\mathcal T=\mathcal T'$。
\end{exercise}
% SUPPLEMENT: id=v14-ch05-ex46-solution; parent=ch05-ex46-initial-seminorm; kind=solution+motivation+hypotheses; origin=independent-derivation
\begin{solution}
\textbf{非满射情形的补写（重构新增）。}$\mathcal T'$ 的每个子基集
\begin{align*}
\{x:\|T_\alpha(x-x_0)\|<\epsilon\}=T_\alpha^{-1}(B(T_\alpha x_0,\epsilon))
\end{align*}
属于 $\mathcal T$，故 $\mathcal T'\subseteq\mathcal T$。反向只须证明 $\mathcal T$ 的每个球原像 $V=T_\alpha^{-1}(B(y_0,\epsilon))$ 属于 $\mathcal T'$。若 $V=\varnothing$，显然成立；否则对每个 $x_0\in V$，令 $\delta=\epsilon-\|T_\alpha x_0-y_0\|>0$。三角不等式给
\begin{align*}
\{x:\|T_\alpha(x-x_0)\|<\delta\}\subseteq V.
\end{align*}
这些 $\mathcal T'$ 开集随 $x_0$ 遍历 $V$ 时覆盖 $V$，所以 $V$ 开，得到反包含。若映射族为空，两边均为平凡拓扑。
\end{solution}
\begin{remark}
\textbf{动机与条件说明（重构新增）。}球中心是否位于映像中不影响其原像的局部描述；只需在每个实际命中点处缩小半径。线性性把两个像的差变成差的像，使这些邻域恰是半范数球。无需单射、满射、完备性或任何事先给定的 $\mathcal X$ 拓扑。
\end{remark}
% END SUPPLEMENT: id=v14-ch05-ex46-solution
% SOURCE p=185; printed=170; object=ch05-ex47-weak-boundedness; number=47; region=middle-lower
\begin{exercise}\label{exercise:ch05-weak-convergent-bounded}
设 $\mathcal X,\mathcal Y$ 为 Banach 空间。
\begin{enumerate}
\item\label{item:ch05-wot-sequence-bounded} 若 $\{T_n\}_{n=1}^{\infty}\subseteq L(\mathcal X,\mathcal Y)$，且 $T_n\to T$ 弱收敛（或强收敛），则 $\sup_n\|T_n\|<\infty$。
\item\label{item:ch05-weak-and-star-sequences-bounded} $\mathcal X$ 中每个弱收敛序列和 $\mathcal X^*$ 中每个弱星收敛序列，都关于范数有界。
\end{enumerate}
\end{exercise}
% SUPPLEMENT: id=v14-ch05-ex47-solution; parent=ch05-ex47-weak-boundedness; kind=solution+motivation+hypotheses; origin=independent-derivation
\begin{solution}
\textbf{两次一致有界原理（重构新增）。}先证明第二问中关于弱收敛的断言，这也适用于任何赋范空间 $Z$。若 $z_n\to z$ 弱收敛，则对每个 $g\in Z^*$，标量列 $g(z_n)$ 收敛，因而有界。将 $Jz_n:g\mapsto g(z_n)$ 看成 Banach 空间 $Z^*$ 上的连续线性泛函，应用一致有界原理，得 $\sup_n\|Jz_n\|<\infty$。自然嵌入等距，故 $\sup_n\|z_n\|<\infty$。

对\subcref{exercise:ch05-weak-convergent-bounded}{item:ch05-wot-sequence-bounded}，固定 $x\in\mathcal X$，弱算子收敛给 $T_nx\to Tx$ 弱收敛；由刚才的结论，$\sup_n\|T_nx\|<\infty$。在 Banach 空间 $\mathcal X$ 上再次对族 $\{T_n\}$ 使用一致有界原理，得到 $\sup_n\|T_n\|<\infty$。强算子收敛蕴含弱算子收敛，故包含在此证明中。

对\subcref{exercise:ch05-weak-convergent-bounded}{item:ch05-weak-and-star-sequences-bounded} 的剩余断言，若 $f_n\to f$ 弱星收敛，则对每个 $x$，$f_n(x)$ 为有界标量序列。在完备空间 $\mathcal X$ 上对 $\{f_n\}$ 用一致有界原理，即得 $\sup_n\|f_n\|<\infty$。
\end{solution}
\begin{remark}
\textbf{动机与逐条件说明（重构新增）。}弱收敛先给每个标量测量的界，再由对偶空间的完备性合并成向量范数界；算子结论还需在输入空间再合并一次。弱收敛序列的有界性不要求原空间完备，因为它的连续对偶自动 Banach；弱星部分和算子部分在本证中使用 $\mathcal X$ 完备，$\mathcal Y$ 完备可省。这里使用“收敛序列整体有界”；一般收敛网的全部取值不必有界，不能照搬到网。
\end{remark}
% END SUPPLEMENT: id=v14-ch05-ex47-solution
% SOURCE p=185; printed=170; object=ch05-ex48-weak-closed-ball; number=48; region=lower
\begin{exercise}\label{exercise:ch05-weak-closures-completeness}
设 $\mathcal X$ 为 Banach 空间。
\begin{enumerate}
\item\label{item:ch05-unit-ball-weakly-closed} 范数闭单位球 $B=\{x\in\mathcal X:\|x\|\le1\}$ 也是弱闭的。（使用\subcref{theorem:ch05-dual-separation-embedding}{item:ch05-canonical-isometry}。）
\item\label{item:ch05-bounded-weak-closure} 若 $E\subseteq\mathcal X$ 关于范数有界，则它的弱闭包也关于范数有界。
\item\label{item:ch05-bounded-weak-star-closure} 若 $F\subseteq\mathcal X^*$ 关于范数有界，则它的弱星闭包也关于范数有界。
\item\label{item:ch05-weak-star-sequential-completeness} $\mathcal X^*$ 中每个弱星 Cauchy 列都收敛。（使用\cref{exercise:ch05-pointwise-operator-limit}。）
\end{enumerate}
\end{exercise}
% SUPPLEMENT: id=v14-ch05-ex48-solution; parent=ch05-ex48-weak-closed-ball; kind=solution+motivation+hypotheses; origin=independent-derivation
\begin{solution}
\textbf{逐项补写（重构新增）。}
\begin{enumerate}
\item 由自然嵌入的等距性，$\|x\|=\sup_{\|f\|\le1}|f(x)|$，所以
\begin{align*}
B=\bigcap_{\|f\|\le1}\{x:|f(x)|\le1\}.
\end{align*}
每个交项是弱连续函数的闭集原像，故弱闭，任意交也闭。
\item 取 $R>0$ 使 $E\subseteq RB$。同样的半空间交表达式说明 $RB$ 弱闭，故 $\overline E^{\mathrm{weak}}\subseteq RB$，得到界。空集也包含在此论证中。
\item 取 $R>0$ 使 $\|f\|\le R$ 对所有 $f\in F$ 成立。集合
\begin{align*}
RB^*=\bigcap_{x\in\mathcal X}\{f\in\mathcal X^*:|f(x)|\le R\|x\|\}
\end{align*}
弱星闭且包含 $F$，所以也包含其弱星闭包，证得有界。
\item 若 $f_n$ 弱星 Cauchy，则对每个 $x$，标量列 $f_n(x)$ 为 Cauchy 列，故有极限 $f(x)$。有限线性组合可过极限，故 $f$ 线性。每个标量列有界，一致有界原理给 $M=\sup_n\|f_n\|<\infty$，所以 $|f(x)|\le M\|x\|$，即 $f\in\mathcal X^*$。按照定义，$f_n\to f$ 弱星收敛。这正是题中所引逐点极限结果在标量目标时的证明。
\end{enumerate}
\end{solution}
\begin{remark}
\textbf{动机与条件说明（重构新增）。}闭范数球可以写成连续线性测量的闭约束之交，因此弱化拓扑后仍闭；有界集合的闭包不能越出这个闭约束。前三问只需赋范空间，第四问用定义域完备性把逐点标量极限提升为连续泛函。闭包采用相应弱拓扑，不能混淆为范数闭包；第四问是序列完备性，并未证明任意弱星 Cauchy 网都有连续的极限。
\end{remark}
% END SUPPLEMENT: id=v14-ch05-ex48-solution
% SOURCE p=185; printed=170; object=ch05-ex49-weak-smallness; number=49; region=lower-bottom
\begin{exercise}\label{exercise:ch05-weak-topology-smallness}
设 $\mathcal X$ 为无限维 Banach 空间。
\begin{enumerate}
\item\label{item:ch05-weak-open-unbounded} $\mathcal X$ 的每个非空弱开集，以及 $\mathcal X^*$ 的每个非空弱星开集，都关于范数无界。
\item\label{item:ch05-weak-bounded-nowhere-dense} $\mathcal X$ 的每个有界子集在弱拓扑下无处稠密，$\mathcal X^*$ 的每个有界子集在弱星拓扑下无处稠密。（使用\subcref{exercise:ch05-weak-closures-completeness}{item:ch05-bounded-weak-closure} 及\subcref{exercise:ch05-weak-closures-completeness}{item:ch05-bounded-weak-star-closure}。）
\item\label{item:ch05-weak-spaces-meager} $\mathcal X$ 在弱拓扑下于自身中疏，$\mathcal X^*$ 在弱星拓扑下于自身中疏。
\item\label{item:ch05-weak-star-not-invariant-metrizable} $\mathcal X^*$ 上的弱星拓扑不由任何平移不变度量定义。（使用\subcref{exercise:ch05-weak-closures-completeness}{item:ch05-weak-star-sequential-completeness}。）
\end{enumerate}
\end{exercise}
% SUPPLEMENT: id=v14-ch05-ex49-solution; parent=ch05-ex49-weak-smallness; kind=solution+motivation+hypotheses; origin=independent-derivation
\begin{solution}
\textbf{逐项补写（重构新增）。}
\begin{enumerate}
\item 对非空弱开集取 $x_0$，其中包含一个由有限个泛函 $f_1,\ldots,f_m$ 及正半径定义的 $x_0$ 的基本邻域。映射 $x\mapsto(f_1(x),\ldots,f_m(x))$ 的核非零，否则无限维空间可线性嵌入 $K^m$。取核中 $v\ne0$，则整条 $x_0+tv$（$t\in\mathbb R$）都在邻域内，且 $\|x_0+tv\|\ge |t|\|v\|-\|x_0\|\to\infty$。

对弱星开集作相同论证，有限测量改为 $f\mapsto(f(x_1),\ldots,f(x_m))$。此处 $\mathcal X^*$ 也无限维：若它有限维，其对偶 $\mathcal X^{**}$ 有限维，而自然嵌入 $J:\mathcal X\to\mathcal X^{**}$ 单射，与 $\mathcal X$ 无限维矛盾。因此有限测量同样有非零共同核，整条仿射直线给无界性。
\item 前题说明有界集在相应弱拓扑中的闭包仍有界。若闭包有非空内部，就包含一个非空开集，而第一问说它无界，矛盾。因此闭包内部为空，集合无处稠密。
\item 每个空间都是其整数半径闭范数球的可数并，每个球由第二问无处稠密，故在所说拓扑中为疏集。
\item 假设弱星拓扑由平移不变度量 $\rho$ 给出。$\rho(f,g)=\rho(f-g,0)$ 说明度量 Cauchy 列恰为弱星 Cauchy 列。前题第四问使每个这样的序列在弱星拓扑、亦即 $\rho$ 拓扑中收敛，所以 $\rho$ 完备。非空完备度量空间由 Baire 定理在自身中非疏，与第三问矛盾。
\end{enumerate}
\end{solution}
\begin{remark}
\textbf{动机与逐条件说明（重构新增）。}有限个线性测量无法限制无限维空间的所有方向，这是弱开集无界的原因；有界闭包与此结合产生无处稠密性。前三问本身只需无限维赋范空间。第四问还使用 Banach 条件所得的弱星序列完备性，以及平移不变性对两种 Cauchy 条件的识别。结论没有否认某个有界子集上的弱星拓扑可以度量化，下一题正讨论这种情形。
\end{remark}
% END SUPPLEMENT: id=v14-ch05-ex49-solution
% SOURCE p=186; printed=171; object=ch05-ex50-separable-dual-ball; number=50; region=top
\begin{exercise}\label{exercise:ch05-separable-dual-ball-metrizable}
若 $\mathcal X$ 为可分赋范线性空间，则 $\mathcal X^*$ 中闭单位球上的弱星拓扑第二可数，从而可度量化。（但请对照\subcref{exercise:ch05-weak-topology-smallness}{item:ch05-weak-star-not-invariant-metrizable}。）
\end{exercise}
% SUPPLEMENT: id=v14-ch05-ex50-solution; parent=ch05-ex50-separable-dual-ball; kind=solution+motivation+hypotheses; origin=independent-derivation
\begin{solution}
\textbf{在单位球上用可数点测试（重构新增）。}取 $\mathcal X$ 的可数稠密集 $\{x_j\}_{j\ge1}$，有限或单点情形可重复编号。记对偶闭单位球为 $B^*$，令 $E(f)=(f(x_1),f(x_2),\ldots)\in K^{\mathbb N}$。若 $E(f)=E(g)$，则连续泛函 $f-g$ 在稠密集上为零，从而处处为零，所以 $E$ 单射。

先证明在 $B^*$ 上这些可数评价已生成全部弱星拓扑。对任意 $x\in\mathcal X$、$f,g\in B^*$，有
\begin{align*}
|(f-g)(x)|\le |(f-g)(x_j)|+2\|x-x_j\|.
\end{align*}
给 $\epsilon>0$，选 $j$ 使 $\|x-x_j\|<\epsilon/4$；只要 $|(f-g)(x_j)|<\epsilon/2$，就有 $|(f-g)(x)|<\epsilon$。对有限多个测试点分别这样选取，说明任何弱星基本邻域含可数评价所定义的基本邻域。反包含直接来自这些评价本身属于全部评价族。因此 $E$ 是 $B^*$ 到其像的拓扑嵌入。

$K$ 有可数基（实数取有理端点区间，复数取有理中心、正有理半径圆盘）；可数乘积的有限坐标基本开集也只有可数多个，所以 $K^{\mathbb N}$ 第二可数。其子空间 $E(B^*)$ 第二可数，故 $B^*$ 第二可数。还可直接写出相容度量
\begin{align*}
\rho(f,g)=\sum_{j=1}^{\infty}2^{-j}\min\{1,|f(x_j)-g(x_j)|\}.
\end{align*}
度量公理与拓扑一致性由\cref{exercise:ch05-prove-seminorm-metrization} 的同一逐项估计得到；上述单射性保证正定。这给出所需度量化，而不需再调用一般度量化定理。
\end{solution}
\begin{remark}
\textbf{动机与逐条件说明（重构新增）。}可分性把所有输入点用可数个点逼近；单位球的统一范数界使点的逼近误差对所有泛函同时受控。缺少这一步，不能把整个对偶空间的弱星拓扑也换成可数评价拓扑。任意固定范数有界子集上同样可行。无需 $\mathcal X$ 完备；与 Alaoglu 定理结合，得到可分赋范空间的对偶单位球为紧度量空间。
\end{remark}
% END SUPPLEMENT: id=v14-ch05-ex50-solution
% SOURCE p=186; printed=171; object=ch05-ex51-weak-closed-subspace; number=51; region=upper
\begin{exercise}\label{exercise:ch05-norm-weak-closed-subspace}
赋范向量空间 $\mathcal X$ 的向量子空间范数闭，当且仅当它弱闭。（不过，$\mathcal X^*$ 的范数闭子空间未必弱星闭，除非 $\mathcal X$ 自反；见\subcref{exercise:ch05-goldstine-reflexive-compactness}{item:ch05-goldstine-density}。）
\end{exercise}
% SUPPLEMENT: id=v14-ch05-ex51-solution; parent=ch05-ex51-weak-closed-subspace; kind=solution+motivation+hypotheses; origin=independent-derivation
\begin{solution}
\textbf{闭子空间的线性分离（重构新增）。}若 $M$ 弱闭，因为弱拓扑包含于范数拓扑，恒等映射从范数拓扑到弱拓扑连续，故 $M$ 范数闭。反之，若 $M$ 范数闭，Hahn--Banach 闭子空间分离给：对每个 $x\notin M$，存在 $f\in\mathcal X^*$ 满足 $f|_M=0$ 且 $f(x)\ne0$。所以
\begin{align*}
M=\bigcap_{f\in M^0}\ker f,
\end{align*}
右边为弱闭集之交，故 $M$ 弱闭。

补证括号中的区别。若 $\mathcal X$ 自反，$\mathcal X^*$ 的弱星与弱拓扑相同，刚证结果适用。若 $\mathcal X$ 不自反，取 $F\in\mathcal X^{**}\setminus J\mathcal X$。它非零，$M=\ker F$ 为范数闭子空间。假如 $M$ 也弱星闭，取 $h$ 使 $F(h)=1$，并取弱星零点基本邻域
\begin{align*}
U=\{u:|u(x_j)|<\epsilon_j, 1\le j\le m\},\qquad (h+U)\cap M=\varnothing.
\end{align*}
若某个 $u\in U$ 满足 $|F(u)|\ge1$，则 $v=-u/F(u)$ 仍在 $U$，而 $F(h+v)=0$，矛盾。因此 $|F(u)|<1$ 于 $U$；对任意 $\delta>0$，$|F|<\delta$ 于开邻域 $\delta U$，故 $F$ 弱星连续。由正文证明的弱星连续对偶等于 $J\mathcal X$，得到矛盾。所以 $M$ 范数闭却不弱星闭。
\end{solution}
\begin{remark}
\textbf{动机与逐条件说明（重构新增）。}子空间的闭性可以用“在它上面为零”的连续线性测量完整描述；任意闭集没有这一个线性核交表示，不能据此推广。范数闭性保证分离点与子空间的正距离及 Hahn--Banach 分离适用，完备性不需要。弱星拓扑只允许原空间评价作为连续线性测量，不自反时会遗漏某些双对偶泛函，所以相应闭超平面可能不再闭。
\end{remark}
% END SUPPLEMENT: id=v14-ch05-ex51-solution
% SOURCE p=186; printed=171; object=ch05-ex52-goldstine; number=52; region=middle
\begin{exercise}\label{exercise:ch05-goldstine-reflexive-compactness}
设 $\mathcal X$ 为 Banach 空间，$f_1,\ldots,f_n$ 为 $\mathcal X^*$ 中线性无关的元素。
\begin{enumerate}
\item\label{item:ch05-finite-evaluation-annihilator} 定义 $T:\mathcal X\to\mathbb C^n$ 为 $Tx=(f_1(x),\ldots,f_n(x))$。若 $\mathcal N=\{x:Tx=0\}$，$\mathcal M$ 为 $f_1,\ldots,f_n$ 的线性张成空间，则按\cref{exercise:ch05-annihilators} 的记号有 $\mathcal M=\mathcal N^0$，从而 $\mathcal M^*$ 同构于 $(\mathcal X/\mathcal N)^*$。
\item\label{item:ch05-finite-evaluation-interpolation} 若 $F\in\mathcal X^{**}$，则对任意 $\epsilon>0$，存在 $x\in\mathcal X$，使 $F(f_j)=f_j(x)$（$j=1,\ldots,n$），且 $\|x\|\le(1+\epsilon)\|F\|$。（$F|_{\mathcal M}$ 可等同于 $(\mathcal X/\mathcal N)^{**}$ 中的一个元素，再因 $\mathcal X/\mathcal N$ 有限维而等同于其中的一个元素。）
\item\label{item:ch05-embedded-weak-topology} 把 $\mathcal X$ 看成 $\mathcal X^{**}$ 的子空间，则 $\mathcal X^{**}$ 上弱星拓扑在 $\mathcal X$ 上诱导的相对拓扑，就是 $\mathcal X$ 的弱拓扑。
\item\label{item:ch05-goldstine-density} 关于 $\mathcal X^{**}$ 上的弱星拓扑，$\mathcal X$ 在 $\mathcal X^{**}$ 中稠密，且 $\mathcal X$ 的闭单位球在 $\mathcal X^{**}$ 的闭单位球中稠密。
\item\label{item:ch05-reflexive-weak-compact-ball} $\mathcal X$ 自反，当且仅当它的闭单位球弱紧。
\end{enumerate}
\end{exercise}
% SUPPLEMENT: id=v14-ch05-ex52-solution; parent=ch05-ex52-goldstine; kind=solution+source-clarification+motivation+hypotheses; origin=independent-derivation
\begin{remark}
\textbf{第一问对偶层次的说明（重构新增）。}原页第一问末句确为 $\mathcal M^*$ 与 $(\mathcal X/\mathcal N)^*$ 同构。二者有限维且维数相等，抽象同构确实存在；但由所引零化子结论直接得到的自然等距识别是 $\mathcal M\cong(\mathcal X/\mathcal N)^*$，进而 $\mathcal M^*\cong(\mathcal X/\mathcal N)^{**}$。第二问的范数估计须使用后一个自然等距识别，不能用一个任意的抽象同构代替。
\end{remark}
\begin{solution}
\textbf{五问完整补写（重构新增）。}记自然等距嵌入为 $J:\mathcal X\to\mathcal X^{**}$。
\begin{enumerate}
\item 首先 $T$ 满射。否则 $T(\mathcal X)$ 是 $\mathbb C^n$ 的真子空间，存在非零线性泛函 $z\mapsto\sum_j a_jz_j$ 在它上面为零（扩展其基即可构造），于是 $\sum_j a_jf_j=0$，与线性无关矛盾。显然 $\mathcal M\subseteq\mathcal N^0$。若 $g\in\mathcal N^0$，令 $\ell(Tx)=g(x)$；因 $g$ 消去 $\ker T$，它良定义。$T$ 满射使 $\ell$ 为 $\mathbb C^n$ 上的线性泛函，所以 $g=\sum_j a_j f_j\in\mathcal M$。故等号成立。

$\mathcal N$ 是有限个连续泛函核的交，闭；商 $Q=\mathcal X/\mathcal N$ 为 $n$ 维赋范空间。由\subcref{exercise:ch05-annihilators}{item:ch05-quotient-dual-annihilator}，映射 $\alpha:Q^*\to\mathcal M$、$\alpha(h)=h\circ\pi$ 为等距同构。因此 $Q^*,\mathcal M,\mathcal M^*$ 都为 $n$ 维；选基可得到原句所述 $\mathcal M^*\cong Q^*$，有限维范数等价保证这是赋范空间同构。但下面使用的是等距的 $\alpha$。
\item 给 $F\in\mathcal X^{**}$，定义 $G\in Q^{**}$ 为 $G(h)=F(\alpha h)$。因 $\alpha$ 等距，$\|G\|=\|F|_{\mathcal M}\|\le\|F\|$。有限维空间自反，其自然映射等距满射，所以存在 $u\in Q$ 使 $G(h)=h(u)$ 对所有 $h\in Q^*$ 成立，且 $\|u\|=\|G\|\le\|F\|$。

若 $F=0$，取 $x=0$。若 $F\ne0$，由商范数的下确界定义选 $x$ 代表 $u$，使
\begin{align*}
\|x\|<\|u\|+\epsilon\|F\|\le(1+\epsilon)\|F\|.
\end{align*}
对 $f_j=\alpha h_j$，有 $f_j(x)=h_j(\pi x)=h_j(u)=G(h_j)=F(f_j)$，证毕。这里不假设商范数在某个代表处取得最小值。
\item 双对偶上的弱星拓扑由 $F\mapsto F(f)$（$f\in\mathcal X^*$）生成。限制到 $J\mathcal X$ 后，这些测试就是 $Jx\mapsto f(x)$，其有限坐标基本邻域恰为 $\mathcal X$ 的弱基本邻域，所以两种拓扑一致。
\item 给任意 $F\in\mathcal X^{**}$ 及其一个弱星基本邻域，它仅涉及有限多个 $g_1,\ldots,g_m\in\mathcal X^*$。在它们的张成空间内选一组基 $f_1,\ldots,f_n$，应用第二问选 $x$，使 $f_j(x)=F(f_j)$。线性性使 $g_l(x)=F(g_l)$ 对所有 $l$ 成立，所以 $Jx$ 在该邻域中。若所有测试为零，取 $x=0$ 即可。这证明 $J\mathcal X$ 稠密。

现在设 $\|F\|\le1$，考虑闭单位球内 $F$ 的相对弱星基本邻域，仍记其有限测试为 $g_l$，允许误差为 $\delta_l>0$。选 $\eta>0$ 足够小，使 $\eta|F(g_l)|/(1+\eta)<\delta_l$ 对全部 $l$ 成立。由第二问和选基步骤，可取 $x$ 使所有测试值精确相等，且 $\|x\|\le1+\eta$。令 $y=x/(1+\eta)$，则 $\|y\|\le1$，并且
\begin{align*}
|Jy(g_l)-F(g_l)|=\frac{\eta}{1+\eta}|F(g_l)|<\delta_l.
\end{align*}
故 $J(B_{\mathcal X})$ 在 $B_{\mathcal X^{**}}$ 中弱星稠密。后一个球弱星闭，因为它由所有闭约束 $|F(f)|\le\|f\|$ 定义；因此其确为前一个球在整个双对偶中的弱星闭包。
\item 若 $\mathcal X$ 自反，$J$ 等距满射，把 $B_{\mathcal X}$ 对应到 $B_{\mathcal X^{**}}$。以 $\mathcal X^*$ 为原空间应用 Alaoglu 定理，后者弱星紧；由第三问的拓扑对应，前者弱紧。

反之，若 $B_{\mathcal X}$ 弱紧，第三问说明 $J(B_{\mathcal X})$ 弱星紧。双对偶弱星拓扑 Hausdorff，故它弱星闭。由第四问，它又在 $B_{\mathcal X^{**}}$ 中稠密，因此两球相等。任意非零 $F\in\mathcal X^{**}$ 先除以 $\|F\|$ 进入单位球，再缩放其原像，即得 $F\in J\mathcal X$；零向量同样有原像。所以 $J$ 满射，$\mathcal X$ 自反。
\end{enumerate}
\end{solution}
\begin{remark}
\textbf{动机与逐条件说明（重构新增）。}有限个独立测量构成有限维商，在该商中双对偶数据能由真实向量实现；商范数的近似最小代表给出 $1+\epsilon$ 的控制。线性无关性用于 $T$ 满射和维数计算，对任意有限测试可先删去依赖关系。允许任意正 $\epsilon$ 避免假定最短代表存在；缩放后只近似保留测试值，恰适合弱星邻域。

第四问通常称为 Goldstine 稠密性。它说明原空间向量足以逼近有限测试，却不意味着在双对偶范数中稠密。第五问把这种稠密性与紧集闭性结合，才升级为满射。以上证明除题设统一的 Banach 背景外实际上只用赋范性；若弱紧单位球条件成立，所得自反性本身还会推出完备性。实空间将 $\mathbb C^n$ 换成 $\mathbb R^n$ 即可。
\end{remark}
% END SUPPLEMENT: id=v14-ch05-ex52-solution
% SOURCE p=186; printed=171; object=ch05-ex53-strong-products; number=53; region=middle-lower
\begin{exercise}\label{exercise:ch05-strong-operator-products}
设 $\mathcal X$ 为 Banach 空间，$\{T_n\},\{S_n\}$ 为 $L(\mathcal X,\mathcal X)$ 中的序列，且 $T_n\to T$、$S_n\to S$ 强收敛。
\begin{enumerate}
\item\label{item:ch05-strong-moving-input} 若 $\{x_n\}\subseteq\mathcal X$ 且 $\|x_n-x\|\to0$，则 $\|T_nx_n-Tx\|\to0$。（使用\subcref{exercise:ch05-weak-convergent-bounded}{item:ch05-wot-sequence-bounded}。）
\item\label{item:ch05-strong-product-convergence} $T_nS_n\to TS$ 强收敛。
\end{enumerate}
\end{exercise}
% SUPPLEMENT: id=v14-ch05-ex53-solution; parent=ch05-ex53-strong-products; kind=solution+motivation+hypotheses; origin=independent-derivation
\begin{solution}
\textbf{逐项补写（重构新增）。}
\begin{enumerate}
\item 由所引结果，$M=\sup_n\|T_n\|<\infty$。于是
\begin{align*}
\|T_nx_n-Tx\|\le M\|x_n-x\|+\|T_nx-Tx\|\longrightarrow0.
\end{align*}
第一项由输入的范数收敛趋零，第二项由固定 $x$ 上的强收敛趋零。
\item 固定任意 $x\in\mathcal X$。令 $y_n=S_nx$、$y=Sx$，强收敛给 $\|y_n-y\|\to0$。第一问应用于 $T_n,y_n$，得到 $T_nS_nx\to TSx$。因为 $x$ 任意，算子乘积强收敛。各乘积有界线性，故确属所讨论的算子空间。
\end{enumerate}
\end{solution}
\begin{remark}
\textbf{动机与条件说明（重构新增）。}把“算子变了”和“输入变了”分成两项，统一算子界控制后者。Banach 性用于由序列强收敛推出统一界；若预先给出此界，相同证明不需要完备性。第二问只断言强收敛序列的乘积收敛，不直接断言乘法在整个强算子拓扑上联合连续；一般网不能无条件使用这里的序列有界性。
\end{remark}
% END SUPPLEMENT: id=v14-ch05-ex53-solution
\end{document}
